\begin{frame}{讲解：执行优化课真正想纠正什么误解}
  \Core{很多人一谈执行优化就想到 kernel，但系统里真正决定收益兑现的往往是整条路径。}
  \begin{itemize}
    \item 算子更快只说明局部时间变短，不说明等待、I/O 或同步是否一起下降。
    \item serving 场景里的动态形状、阶段切换和新请求插入，会让“训练里成立的最优执行”不再稳定。
    \item 所以这一讲真正要训练的是路径分解能力，而不是只会比较算子数字。
  \end{itemize}
\end{frame}

\begin{frame}{讲解：如何判断一次执行优化是否真的成立}
  \Core{执行优化必须同时给出局部证据、路径解释和端到端结果，缺一不可。}
  \begin{itemize}
    \item 局部证据说明改了什么，例如 I/O 重叠、编译命中、launch 减少。
    \item 路径解释说明为什么这些局部变化会传递到请求级时延或吞吐。
    \item 端到端结果说明收益有没有真的被用户感知到，而不是停留在 profiler 里。
  \end{itemize}
\end{frame}

\begin{frame}{例子：为什么“优化执行路径”不等于“找更快 kernel”（1）}
  \Core{执行优化课最需要纠正的，是把所有性能改进都理解成单点算子提速。}
\begin{block}{算子层优化}
    可以降低单轮执行时间，但未必解决主路径上的等待、搬运和重组代价。
  \end{block}
\end{frame}

\begin{frame}{例子：为什么“优化执行路径”不等于“找更快 kernel”（2）}
  \Core{执行优化课最需要纠正的，是把所有性能改进都理解成单点算子提速。}
\begin{block}{执行路径优化}
    关注的是数据何时准备好、图何时稳定、I/O 是否与计算重叠、资源是否被正确放置。
  \end{block}
\end{frame}

\begin{frame}{例子：为什么“优化执行路径”不等于“找更快 kernel”（3）}
  \Core{执行优化课最需要纠正的，是把所有性能改进都理解成单点算子提速。}
\begin{block}{课程结论}
    真正的执行优化，往往是路径组织与资源编排问题，而不只是更快的算子实现。
  \end{block}
\end{frame}

\begin{frame}{例子：图执行为什么在 serving 里更难（1）}
  \Core{训练里稳定的图，不代表 serving 里也稳定，因为请求长度、batch 组成和阶段切换都在变化。}
\begin{block}{训练场景}
    形状相对稳定，图编译和缓存更容易复用。
  \end{block}
\end{frame}

\begin{frame}{例子：图执行为什么在 serving 里更难（2）}
  \Core{训练里稳定的图，不代表 serving 里也稳定，因为请求长度、batch 组成和阶段切换都在变化。}
\begin{block}{serving 场景}
    prompt 长度、decode 轮数和新请求插入都会破坏图的稳定性。
  \end{block}
\end{frame}

\begin{frame}{例子：图执行为什么在 serving 里更难（3）}
  \Core{训练里稳定的图，不代表 serving 里也稳定，因为请求长度、batch 组成和阶段切换都在变化。}
\begin{block}{教学重点}
    所谓 compile stabilization，本质是在给动态服务路径创造“足够稳定”的形状类。
  \end{block}
\end{frame}

\begin{frame}{代码例子：最小编译开关为什么值得上课讲（1）}
  \Core{即使课堂不深入编译器实现，也应该让学生看到“执行路径优化”如何落到运行时开关。}
\begin{block}{代码片段}
    {\ttfamily\small
    \detokenize{llm = LLM(path, enforce_eager=False, max_model_len=4096)}\par
    \detokenize{// eager / compiled path changes runtime behavior}\par
    }
  \end{block}
\end{frame}

\begin{frame}{代码例子：最小编译开关为什么值得上课讲（2）}
  \Core{即使课堂不深入编译器实现，也应该让学生看到“执行路径优化”如何落到运行时开关。}
\begin{block}{课堂应追问}
    为什么同一个模型在 eager 与 compiled path 下，收益边界和失败模式都可能不同？
  \end{block}
\end{frame}

\begin{frame}{代码例子：最小编译开关为什么值得上课讲（3）}
  \Core{即使课堂不深入编译器实现，也应该让学生看到“执行路径优化”如何落到运行时开关。}
\begin{block}{系统含义}
    执行路径优化必须伴随更细的可验证性要求，而不能只看一次跑得快不快。
  \end{block}
\end{frame}

\begin{frame}{代码例子：双缓冲与 I/O 重叠的伪代码（1）}
  \Core{课程中可以用小段伪代码让学生理解“放置与重叠”比数学公式更接近系统实现。}
\begin{block}{伪代码}
    {\ttfamily\small
    \detokenize{prefetch(next_weight_chunk)}\par
    \detokenize{compute(current_chunk, current_kv)}\par
    \detokenize{synchronize_if_needed()}\par
    }
  \end{block}
\end{frame}

\begin{frame}{代码例子：双缓冲与 I/O 重叠的伪代码（2）}
  \Core{课程中可以用小段伪代码让学生理解“放置与重叠”比数学公式更接近系统实现。}
\begin{block}{课堂应追问}
    这里真正要隐藏的是什么：权重搬运、KV 访问，还是 launch 开销？
  \end{block}
\end{frame}

\begin{frame}{代码例子：双缓冲与 I/O 重叠的伪代码（3）}
  \Core{课程中可以用小段伪代码让学生理解“放置与重叠”比数学公式更接近系统实现。}
\begin{block}{课程结论}
    所谓 I/O-aware，核心并不是“搬得快”，而是“搬运时尽量不阻塞主计算路径”。
  \end{block}
\end{frame}

\begin{frame}{例子：speculative decoding 为什么也是执行路径问题（1）}
  \Core{speculative decoding 不只是算法策略，它还要求 draft 与 target 两条执行链在系统里协同。}
\begin{block}{理想收益}
    通过先草拟、再验证，减少 target model 的有效工作量。
  \end{block}
\end{frame}

\begin{frame}{例子：speculative decoding 为什么也是执行路径问题（2）}
  \Core{speculative decoding 不只是算法策略，它还要求 draft 与 target 两条执行链在系统里协同。}
\begin{block}{真实代价}
    需要维护额外模型状态、验证回退逻辑和两条路径之间的同步成本。
  \end{block}
\end{frame}

\begin{frame}{例子：speculative decoding 为什么也是执行路径问题（3）}
  \Core{speculative decoding 不只是算法策略，它还要求 draft 与 target 两条执行链在系统里协同。}
\begin{block}{教学意义}
    这类机制天然属于“执行路径组织”而不只是“模型技巧”。
  \end{block}
\end{frame}

\begin{frame}{研究案例：MLP Materialization 的负结果为什么有价值}
  \Core{profiler-driven optimization 的第一课：不要只优化一个看起来显眼、但不在主导路径上的算子。}
  \begin{itemize}
    \item 实测 serving path 中反复出现 \texttt{\detokenize{MatMulV2 -> SwiGlu -> MatMulV2}}。
    \item profiler 显示 \texttt{\detokenize{SwiGlu}} 约占 1.08\% device op time，而 \texttt{\detokenize{MatMulV2}} 约占 86.72\%。
    \item naive 单 op / inplace 优化难以带来端到端收益，下一步应看 MLP boundary、materialization 与 writeback。
  \end{itemize}
\end{frame}

\begin{frame}{研究案例：执行路径优化要问三个问题}
  \Core{一个优化没有兑现时，课程里要训练学生沿路径定位，而不是简单说“实现不够好”。}
  \begin{enumerate}
    \item 它减少的是 kernel time、launch、I/O、状态搬运，还是等待？
    \item 这个局部变化是否落在 TTFT、TPOT 或 throughput 的主导路径上？
    \item profiler、trace 和端到端结果是否互相解释，还是互相矛盾？
  \end{enumerate}
\end{frame}

\begin{frame}{课堂练习：用路径语言解释一次优化（1）}
  \Core{执行优化课最重要的能力，是把“优化点”翻译成“路径改变”。}
\begin{block}{练习一}
    任选一种优化，说明它减少的是 launch、I/O 等待、图不稳定，还是状态搬运。
  \end{block}
\end{frame}

\begin{frame}{课堂练习：用路径语言解释一次优化（2）}
  \Core{执行优化课最重要的能力，是把“优化点”翻译成“路径改变”。}
\begin{block}{练习二}
    说明为什么“更快 kernel”不能自动推出“更低端到端时延”。
  \end{block}
\end{frame}

\begin{frame}{课堂练习：用路径语言解释一次优化（3）}
  \Core{执行优化课最重要的能力，是把“优化点”翻译成“路径改变”。}
\begin{block}{练习三}
    用一句话回答：这讲里“执行路径”真正指的是什么？
  \end{block}
\end{frame}
